%%%%%%%%%%%%%%%%%%%%%%%%%%%%%%%%%%%%%%%%%%%%%%%%%%%%%%%%%%%%%%%%%%%%%%%%
%                                                                      %
%     File: Thesis_Abstract.tex                                        %
%     Tex Master: Thesis.tex                                           %
%                                                                      %
%     Author: Andre C. Marta                                           %
%     Last modified :  4 Mar 2024                                      %
%                                                                      %
%%%%%%%%%%%%%%%%%%%%%%%%%%%%%%%%%%%%%%%%%%%%%%%%%%%%%%%%%%%%%%%%%%%%%%%%

\section*{Abstract}
\addcontentsline{toc}{section}{Abstract}

Neural sequence generators are typically decoded by selecting the single most probable output under the model, utilizing greedy or beam search as approximations to Maximum A Posteriori (MAP) decoding. Across generative tasks, this likelihood-first choice is frequently misaligned with what human evaluators and standard metrics actually reward. The most confident transcription may contain avoidable semantic errors, and the highest-probability image caption can lack specific visual fidelity compared to alternative hypotheses generated by the same model. This work addresses that fundamental mismatch by adopting Quality-Aware Decoding (QAD), a generate-then-choose paradigm in which multiple diverse candidates are produced and the final output is selected using task-aligned utilities and Minimum Bayes Risk (MBR) reranking.

By adapting QAD to automatic speech recognition (ASR) and image captioning, this thesis systematically evaluates the tension between candidate diversity, acoustic and visual grounding, and inference latency. Empirical experiments on the LibriSpeech and MS COCO benchmarks using Whisper and BLIP architectures reveal critical vulnerabilities in contemporary referenceless evaluation. Specifically, the results expose a severe hallucination trap in ASR, demonstrating that referenceless text-centric estimators heavily reward fluent but entirely fabricated transcriptions generated by Diverse Beam Search. In image captioning, the data illustrates a profound modality gap, proving that optimizing strictly for reference-free visual alignment via CLIPScore actively degrades structural stylometry and consensus-based metrics like CIDEr. To resolve these bottlenecks, a Two-Stage MBR framework is proposed and validated. By aggressively pruning candidate pools with fast quality estimators before applying rigorous reference-based consensus, the two-stage pipeline successfully restores multimodal fidelity, mitigates the hallucination trap, and drastically reduces the quadratic computational cost of traditional MBR, establishing a highly efficient paradigm for generative deployment.

\vfill
\textbf{\Large Keywords:} Quality-Aware Decoding, Automatic Speech Recognition, Image Captioning, Minimum Bayes Risk, Quality Estimation, Hallucination
